% ============================================================
% INICIO DEL PAPER: TÍTULO, AUTORES Y AFILIACIONES
% ============================================================

\begin{frontmatter}

\title{
Modelación numérica de la distribución térmica en una placa electrónica idealizada mediante Diferencias Finitas, Gauss--Seidel y Gradiente Conjugado
}

% Verificar el título definitivo, los nombres de los autores y la afiliación.
\author[utec]{Leon Vasquez, Jose Roberto}
\author[utec]{Contreras Sosa, Hector Fabian}
\author[utec]{Rosales Esteban, Antony Yonel}
\author[utec]{Oceda Chavez, Elizabeth Emperatriz}
\author[utec]{Roman Ramirez, Andre Moises}

\affiliation[utec]{
    organization={Universidad de Ingeniería y Tecnología (UTEC)},
    city={Lima},
    country={Perú}
}

\abstracttitle{Resumen}
\keywordtitle{Palabras clave}

\begin{abstract}
La transferencia de calor en componentes electr\'onicos produce campos de temperatura no uniformes cuya distribuci\'on puede estudiarse mediante modelos num\'ericos. Este trabajo modela la conducci\'on estacionaria en una placa electr\'onica idealizada con fuentes volum\'etricas localizadas asociadas con la CPU y la GPU. El objetivo es analizar la respuesta t\'ermica y comparar los m\'etodos de Gauss--Seidel y Gradiente Conjugado al resolver el sistema discretizado. Se aplic\'o el M\'etodo de Diferencias Finitas con aproximaciones centrales, condiciones de Dirichlet y una condici\'on inicial uniforme de $25\ ^\circ$C. En la malla base de $49\times33$ nodos, la temperatura m\'axima calculada fue $77.7592\ ^\circ$C; con tolerancia de residuo relativo de $10^{-8}$, Gauss--Seidel convergi\'o en 2145 iteraciones y Gradiente Conjugado en 127. Tambi\'en se evaluaron el refinamiento espacial y la sensibilidad ante cambios de potencia de las fuentes. Las diferencias entre mallas disminuyeron al refinar la discretizaci\'on, mientras que Gradiente Conjugado requiri\'o menos iteraciones en el escenario base. Los valores describen el modelo idealizado y no son mediciones de una laptop real.
\end{abstract}

\begin{keyword}
Conducci\'on t\'ermica; Diferencias finitas; Gauss--Seidel; Gradiente conjugado; Convergencia num\'erica
\end{keyword}

\end{frontmatter}
